% Generated draft: requires mathematical and source review.
\documentclass[11pt]{article}
\usepackage[margin=25mm]{geometry}
\usepackage{fontspec}
\setmainfont{Latin Modern Roman}
\usepackage{amsmath,amssymb,mathtools,unicode-math}
\setmathfont{Latin Modern Math}
\usepackage{enumitem,needspace,xurl,hyperref}
\setlength{\parindent}{0pt}
\setlength{\parskip}{5pt}
\setlength{\emergencystretch}{3em}
\newcommand{\N}{\mathbb{N}}
\newcommand{\Z}{\mathbb{Z}}
\newcommand{\Q}{\mathbb{Q}}
\newcommand{\R}{\mathbb{R}}
\newcommand{\C}{\mathbb{C}}
\newcommand{\F}{\mathbb{F}}
\newcommand{\A}{\mathbb{A}}
\newcommand{\PP}{\mathbb{P}}
\newcommand{\E}{\mathbb{E}}
\newcommand{\eps}{\varepsilon}
\newcommand{\dd}{\,\mathrm{d}}
\begin{document}
\section*{UnsolvedMath ID: 30001560}
\textbf{Source problem code:} catalogue entry
\subsection{Definitions and mathematical statement}
Let \(I=\{0,1,2,\ldots\}\). For an integer sequence \(b=(b_k)_{k\in I}\), a starting position \(s\in I\), and a positive integer \(n\), define the Hankel determinant
\[
H_n^{(s)}(b)=\det\bigl(b_{s+i+j}\bigr)_{0\leq i,j<n}\in\Z.
\]
The questions below interpret ``every Hankel determinant'' as quantifying over every starting position and every positive order. The supplied assessment mentions all positions, but the statement does not explicitly define its Hankel indexing convention; this interpretation requires review against the original source.

Let \(\varphi\) be the morphism on finite words over \(\{1,2,3,4\}\), extended by concatenation from
\[
\varphi(1)=12,\qquad \varphi(2)=23,\qquad
\varphi(3)=14,\qquad \varphi(4)=32.
\]
Since \(\varphi(1)\) begins with \(1\), its iterates determine the infinite word
\[
a=\varphi^\omega(1)=\lim_{r\to\infty}\varphi^r(1),
\]
where the limit means that each finite prefix eventually stabilizes. Regard its letters as integers and index them as \(a=(a_k)_{k\in I}\).

\textbf{Question 1.} Is it true that
\[
\forall s\in I\;\forall n\geq1,\qquad H_n^{(s)}(a)\neq0
\quad\text{in }\Z?
\]

\textbf{Question 2.} The source introduces the following question with ``More generally'': is it true that
\[
\begin{gathered}
\forall m\in\Z,\ m\geq2,\quad
\forall A\subseteq\Z\text{ finite and nonempty},\quad
\forall b\in A^I,\\
\exists s\in I\;\exists n\geq1,\qquad
H_n^{(s)}(b)\equiv0\pmod m?
\end{gathered}
\]
Equivalently, for each integer \(m\geq2\), no infinite sequence taking values in a finite subset of \(\Z\) has all its Hankel determinants nonzero modulo \(m\). Here nonzero modulo \(m\) means not congruent to zero; it does not require invertibility modulo a composite \(m\). Question 1 concerns nonvanishing in \(\Z\), whereas Question 2 concerns modular nonvanishing; no implication between the two formulations is asserted here.
\subsection{Short English statement}
The first question asks whether every positive-order Hankel determinant at every starting position of the specified morphic sequence is a nonzero integer. The second asks whether, for each modulus \(m\geq2\), every integer sequence with finite range has some Hankel determinant divisible by \(m\). The all-position indexing convention requires source review. Both are supplied as open questions, not independently verified claims.
\subsection{Sources}
\begin{itemize}
\item Local UnsolvedMath registry, permanent ID \(30001560\), title: \emph{Nonvanishing Hankel Determinants of Morphic Sequences}; exact source problem code: \texttt{catalogue entry}. The supplied record belongs to the Oberwolfach Reports Open Problems collection and lists proposed year \(2010\); no proposer is identified.
\item Supplied original workshop reference: \emph{Mini-Workshop: Combinatorics on Words (2011)}, Oberwolfach Reports, pp.~2195--2244, DOI \texttt{10.4171/owr/2010/37}. \url{https://doi.org/10.4171/owr/2010/37}. The report is described as documenting workshop discussions and their open problems. The supplied title year \(2011\) and proposed year \(2010\) are retained without reconciliation.
\item Supplied literature source, labeled ``current open-problems survey'': \url{https://cs.uwaterloo.ca/~shallit/Talks/bc4.pdf}. No author, title, or publication date for this survey is supplied.
\item Literature assessment dated exactly \texttt{2026-08-21}: the supplied curation labels the problem open and reports that the specified morphism and its all-position nonzero-Hankel-determinant property remain listed as open. It explicitly cautions that finite computation does not constitute a proof. This is supplied dated curation derived from the OWR record, not an independently verified literature assessment; time-sensitive status and cited claims require independent verification.
\end{itemize}
\end{document}
